\chapter{Core Data Structures: The Classic Network}
\label{ch:data}

\glance{The classic ABC design is an \id{Abc\_Ntk\_t} --- a flat array of
\id{Abc\_Obj\_t} objects (PIs, POs, latches, nodes, \dots) connected by
\emph{integer-ID} fanin/fanout edges resolved through \id{pNtk->vObjs}. A
network has a \emph{type} (netlist / logic / strash) and a \emph{functionality}
(SOP / BDD / AIG / map). The most important type is \textbf{STRASH}: a
structurally hashed AIG of two-input AND gates with complemented edges, built by
\cmd{strash} and required by \cmd{rewrite}, \cmd{refactor}, and \cmd{resyn2}.
All definitions in this chapter live in \file{src/base/abc/abc.h} and its
implementation files \file{abcNtk.c}, \file{abcObj.c}, and \file{abcAig.c}.}

\section{The two central structs}

Every design object is an \id{Abc\_Obj\_t} (declared as \id{struct Abc\_Obj\_t\_}
in \file{src/base/abc/abc.h}). It is deliberately small (48/72 bytes on
32/64-bit builds) because a real design holds millions of them.

\begin{lstlisting}[style=abccode]
struct Abc_Obj_t_     // 48/72 bytes (32-bits/64-bits)
{
    Abc_Ntk_t *   pNtk;          // the host network
    Abc_Obj_t *   pNext;         // the next pointer in the hash table
    int           Id;            // the object ID
    unsigned      Type    :  4;  // the object type
    unsigned      fMarkA  :  1;  // the multipurpose mark
    unsigned      fMarkB  :  1;
    unsigned      fMarkC  :  1;
    unsigned      fPhase  :  1;  // phase of equivalent node
    unsigned      fExor   :  1;  // AIG node is a root of EXOR
    unsigned      fPersist:  1;  // persistant AIG node
    unsigned      fCompl0 :  1;  // compl. attr. of the first fanin
    unsigned      fCompl1 :  1;  // compl. attr. of the second fanin
    unsigned      Level   : 20;  // the level of the node
    Vec_Int_t     vFanins;       // the array of fanins  (object IDs)
    Vec_Int_t     vFanouts;      // the array of fanouts (object IDs)
    void *        pDataComp;
    union { void * pData; int iData; };  // network-specific data
    union { void * pTemp; Abc_Obj_t * pCopy; int iTemp; float dTemp; };
};
\end{lstlisting}

The network itself, \id{struct Abc\_Ntk\_t\_}, owns the object arrays, the
per-type counts, the name manager, and a \emph{functionality manager}
(\id{pManFunc}) whose concrete type depends on \id{ntkFunc} (an AIG manager, a
BDD manager, or a SOP memory manager).

\begin{lstlisting}[style=abccode]
struct Abc_Ntk_t_ {
    Abc_NtkType_t ntkType;       // netlist / logic / strash
    Abc_NtkFunc_t ntkFunc;       // sop / bdd / aig / map / ...
    char *        pName;
    Nm_Man_t *    pManName;      // stores object names
    Vec_Ptr_t *   vObjs;         // ALL objects; index == Id
    Vec_Ptr_t *   vPis, * vPos;  // primary inputs / outputs
    Vec_Ptr_t *   vCis, * vCos;  // combinational in (PI,BO) / out (PO,BI)
    Vec_Ptr_t *   vBoxes;        // latches and white/black boxes
    int  nObjCounts[ABC_OBJ_NUMBER]; // live objects by type
    int  nObjs;                  // number of live objects
    Mem_Fixed_t * pMmObj;        // memory manager for objects
    void *        pManFunc;      // AIG / BDD / SOP manager
    /* ... timing, cuts, EXDC, attributes, traversal IDs ... */
};
\end{lstlisting}

\section{Network type and functionality}

\id{ntkType} and \id{ntkFunc} are orthogonal enums. Not every combination is
legal; the header documents the supported grid (a strashed network is always an
AIG, for instance).

\begin{center}
\begin{tabularx}{\linewidth}{@{}l l X@{}}
\toprule
\textbf{Constant} & \textbf{Value} & \textbf{Meaning} \\
\midrule
\id{ABC\_NTK\_NONE}    & 0 & unknown \\
\id{ABC\_NTK\_NETLIST} & 1 & PIs/POs, latches, nodes, \emph{and nets} \\
\id{ABC\_NTK\_LOGIC}   & 2 & PIs/POs, latches, nodes (no nets) \\
\id{ABC\_NTK\_STRASH}  & 3 & structurally hashed AIG (2-input ANDs, c-attributes on edges) \\
\id{ABC\_NTK\_OTHER}   & 4 & unused \\
\bottomrule
\end{tabularx}
\end{center}

\begin{center}
\begin{tabularx}{\linewidth}{@{}l l X@{}}
\toprule
\textbf{Constant} & \textbf{Value} & \textbf{Node functionality} \\
\midrule
\id{ABC\_FUNC\_SOP}      & 1 & sum-of-products (a string per node) \\
\id{ABC\_FUNC\_BDD}      & 2 & binary decision diagrams (CUDD) \\
\id{ABC\_FUNC\_AIG}      & 3 & and-inverter graphs \\
\id{ABC\_FUNC\_MAP}      & 4 & standard-cell library gates \\
\id{ABC\_FUNC\_BLIFMV}   & 5 & BLIF-MV node functions \\
\id{ABC\_FUNC\_BLACKBOX} & 6 & black box (contents unknown) \\
\bottomrule
\end{tabularx}
\end{center}

The header supplies inline predicates for both axes, e.g.\ \id{Abc\_NtkIsStrash},
\id{Abc\_NtkIsLogic}, \id{Abc\_NtkHasAig}, and combined ones like
\id{Abc\_NtkIsSopLogic}. The value \id{ABC\_NTK\_STRASH} is the gateway to
optimization: commands such as \cmd{rewrite} and \cmd{resyn2} assert an AIG, so a
logic network must be converted with \cmd{strash} first.

\section{Object types}

Each object carries a 4-bit \id{Type} field drawn from \id{Abc\_ObjType\_t}.
Predicates like \id{Abc\_ObjIsNode}, \id{Abc\_ObjIsCi}, \id{Abc\_ObjIsCo}, and
\id{Abc\_ObjIsBox} classify objects; \id{Abc\_ObjIsCi} is true for PIs and box
outputs, \id{Abc\_ObjIsCo} for POs and box inputs.

\begin{center}
\begin{tabularx}{\linewidth}{@{}l l X@{}}
\toprule
\textbf{Constant} & \textbf{Val} & \textbf{Role} \\
\midrule
\id{ABC\_OBJ\_CONST1}   & 1  & constant-1 node (AIG only; always Id 0's peer) \\
\id{ABC\_OBJ\_PI}       & 2  & primary input terminal \\
\id{ABC\_OBJ\_PO}       & 3  & primary output terminal \\
\id{ABC\_OBJ\_BI}       & 4  & box input terminal (a CO) \\
\id{ABC\_OBJ\_BO}       & 5  & box output terminal (a CI) \\
\id{ABC\_OBJ\_NET}      & 6  & net (netlist networks only) \\
\id{ABC\_OBJ\_NODE}     & 7  & logic node / AIG AND gate \\
\id{ABC\_OBJ\_LATCH}    & 8  & latch (a box) \\
\id{ABC\_OBJ\_WHITEBOX} & 9  & box with known contents \\
\id{ABC\_OBJ\_BLACKBOX} & 10 & box with unknown contents \\
\bottomrule
\end{tabularx}
\end{center}

Latch initial values are a separate enum \id{Abc\_InitType\_t}
(\id{ABC\_INIT\_ZERO}, \id{ABC\_INIT\_ONE}, \id{ABC\_INIT\_DC}), stored directly
in the latch's \id{pData} field via helpers such as \id{Abc\_LatchSetInit0}.

\section{Edges are integer IDs, not pointers}

A fanin/fanout is stored as an \emph{integer object ID} inside the
\id{vFanins}/\id{vFanouts} vectors, and resolved lazily against
\id{pNtk->vObjs}. The inline accessors make this indirection invisible:

\begin{lstlisting}[style=abccode]
static inline Abc_Obj_t * Abc_ObjFanin( Abc_Obj_t * pObj, int i )
    { return (Abc_Obj_t *)pObj->pNtk->vObjs->pArray[ pObj->vFanins.pArray[i] ]; }
static inline int Abc_ObjFaninNum ( Abc_Obj_t * pObj ) { return pObj->vFanins.nSize;  }
static inline int Abc_ObjFanoutNum( Abc_Obj_t * pObj ) { return pObj->vFanouts.nSize; }
\end{lstlisting}

Because edges are IDs, an object can be deleted (its \id{vObjs} slot set to
\id{NULL}) without dangling pointers elsewhere; iteration macros skip the holes.
The \id{Vec\_*} vector types used here are covered in the utilities chapter
(\ref{ch:util}) and not detailed further.

\subsection{Iteration macros}

Traversal is done almost exclusively through macros so that \id{NULL} holes and
type filters are handled uniformly. The essential ones:

\begin{center}
\begin{tabularx}{\linewidth}{@{}l X@{}}
\toprule
\textbf{Macro} & \textbf{Iterates over} \\
\midrule
\id{Abc\_NtkForEachObj(pNtk,pObj,i)}      & every live object in \id{vObjs} \\
\id{Abc\_NtkForEachNode(pNtk,pNode,i)}    & objects with \id{Type == NODE} \\
\id{Abc\_AigForEachAnd(pNtk,pNode,i)}     & two-input AND gates of a strashed AIG \\
\id{Abc\_NtkForEachPi / \dots ForEachPo}  & primary inputs / outputs \\
\id{Abc\_NtkForEachCi / \dots ForEachCo}  & combinational inputs / outputs \\
\id{Abc\_NtkForEachLatch(pNtk,pObj,i)}    & latches (walks \id{vBoxes}) \\
\id{Abc\_ObjForEachFanin(pObj,pFanin,i)}  & fanins of one object \\
\id{Abc\_ObjForEachFanout(pObj,pFanout,i)}& fanouts of one object \\
\bottomrule
\end{tabularx}
\end{center}

\section{Complemented edges: two encodings}

ABC represents an inverter as an \emph{attribute on an edge}, never as a
standalone node. There are two distinct encodings for the same idea.

\paragraph{Stored bits (structural AIG edges).} An AND gate's two incoming edges
carry complement bits \id{fCompl0} and \id{fCompl1} in the object itself. The
combination of a fanin and its bit is the \emph{child}:

\begin{lstlisting}[style=abccode]
static inline int Abc_ObjFaninC0( Abc_Obj_t * p ) { return p->fCompl0; }
static inline Abc_Obj_t * Abc_ObjChild0( Abc_Obj_t * p )
    { return Abc_ObjNotCond( Abc_ObjFanin0(p), Abc_ObjFaninC0(p) ); }
\end{lstlisting}

\paragraph{Pointer tagging (in-flight edges).} While building or querying the
AIG, an edge is passed around as an \id{Abc\_Obj\_t*} whose least-significant bit
encodes complementation. Three one-line inlines manage the tag:

\begin{lstlisting}[style=abccode]
static inline int Abc_ObjIsComplement( Abc_Obj_t * p )
    { return (int)((ABC_PTRUINT_T)p & (ABC_PTRUINT_T)01); }
static inline Abc_Obj_t * Abc_ObjRegular( Abc_Obj_t * p )
    { return (Abc_Obj_t *)((ABC_PTRUINT_T)p & ~(ABC_PTRUINT_T)01); }
static inline Abc_Obj_t * Abc_ObjNot( Abc_Obj_t * p )
    { return (Abc_Obj_t *)((ABC_PTRUINT_T)p ^  (ABC_PTRUINT_T)01); }
\end{lstlisting}

\id{Abc\_ObjNotCond(p,c)} flips the bit only when \id{c} is nonzero. Because the
LSB is stolen, objects must be at least 2-byte aligned (they are). The tagged
form is what \id{Abc\_AigAnd} and friends accept and return; the stored-bit form
is what persists once a node is created.

\section{Structural hashing (the AIG)}

The strash manager \id{struct Abc\_Aig\_t\_} (in \file{src/base/abc/abcAig.c})
is the \id{pManFunc} of a strashed network. Its heart is a chained hash table of
AND gates:

\begin{lstlisting}[style=abccode]
struct Abc_Aig_t_ {
    Abc_Ntk_t *  pNtkAig;    // the AIG network
    Abc_Obj_t *  pConst1;    // the constant 1 object (not a node!)
    Abc_Obj_t ** pBins;      // the table bins (chained by pObj->pNext)
    int          nBins;      // size of the table
    int          nEntries;   // total entries
    /* ... incremental-update scratch vectors ... */
};
\end{lstlisting}

The invariant, quoted from the file header, is that \emph{for each AND gate
there is no other AND gate with the same children}, constants are propagated,
and there are no single-input (buffer/inverter) nodes. Two AND gates are
``the same'' when they have identical (fanin, complement) pairs. To make the key
order-independent, the arguments are sorted by regular-object \id{Id} before
hashing:

\begin{lstlisting}[style=abccode]
Abc_Obj_t * Abc_AigAnd( Abc_Aig_t * pMan, Abc_Obj_t * p0, Abc_Obj_t * p1 )
{
    Abc_Obj_t * pAnd;
    if ( (pAnd = Abc_AigAndLookup( pMan, p0, p1 )) )  // reuse existing node
        return pAnd;
    return Abc_AigAndCreate( pMan, p0, p1 );          // else make a new one
}
\end{lstlisting}

\id{Abc\_AigAndLookup} first collapses the trivial cases --- $x\wedge x = x$,
$x\wedge\overline{x}=0$, and constant folding against \id{pConst1} --- then
computes \id{Abc\_HashKey2} (which mixes both fanin IDs and both complement
bits) and walks the bin's \id{pNext} chain looking for a child-for-child match.
\id{Abc\_AigAndCreate} orders the arguments, allocates the node, sets its level
as $1+\max(\text{fanin levels})$, marks EXOR/phase status, and links it into the
front of its bin:

\begin{lstlisting}[style=abccode]
if ( Abc_ObjRegular(p0)->Id > Abc_ObjRegular(p1)->Id )  // canonical fanin order
    pAnd = p0, p0 = p1, p1 = pAnd;
pAnd = Abc_NtkCreateNode( pMan->pNtkAig );
Abc_ObjAddFanin( pAnd, p0 );  Abc_ObjAddFanin( pAnd, p1 );
pAnd->Level = 1 + Abc_MaxInt( Abc_ObjRegular(p0)->Level, Abc_ObjRegular(p1)->Level );
Key = Abc_HashKey2( p0, p1, pMan->nBins );
pAnd->pNext = pMan->pBins[Key];  pMan->pBins[Key] = pAnd;  pMan->nEntries++;
\end{lstlisting}

Because identical ANDs are shared rather than duplicated, every rewrite that
recreates a subgraph automatically deduplicates it. This is the source of the
node-count ``gain'' reported by optimization passes: re-hashing a sloppy AIG (for
example during \id{Abc\_NtkDup}) can collapse it to fewer nodes, and ABC prints a
warning when duplication alone reduces the node count. Boolean helpers
\id{Abc\_AigOr}, \id{Abc\_AigXor}, and \id{Abc\_AigMux} are all expressed in
terms of \id{Abc\_AigAnd} and pointer complementation.

\section{Scratch fields for dataflow passes}

The last two struct members exist purely so that passes can stash temporary
per-object data without allocating side tables. They are unions, so a pass picks
one interpretation at a time.

\begin{center}
\begin{tabularx}{\linewidth}{@{}l l X@{}}
\toprule
\textbf{Field} & \textbf{Union} & \textbf{Typical use} \\
\midrule
\id{pCopy}   & with \id{pTemp/iTemp/dTemp} & maps an old object to its image in a new network \\
\id{pData}   & with \id{iData}             & the node's function (SOP string, BDD, gate, latch init, equiv class) \\
\id{iTemp}\,/\,\id{dTemp} & with \id{pTemp} & scratch int/float: simulation values, SAT variables \\
\id{fMarkA/B/C} & bitfield                & visited flags for local traversals \\
\id{Level}   & 20-bit field                & structural depth from the CIs \\
\bottomrule
\end{tabularx}
\end{center}

\id{pCopy} is the workhorse of every ``build a fresh network'' pass: helpers
such as \id{Abc\_ObjChild0Copy} read \id{Abc\_ObjFanin0(pObj)->pCopy} and
re-apply the stored complement bit, which is exactly how \id{Abc\_NtkDup}
reconstructs edges in the new network. Traversal IDs (\id{nTravIds},
\id{vTravIds}) provide an $O(1)$ ``visited this pass'' test via
\id{Abc\_NodeIsTravIdCurrent} without clearing marks between runs. Node levels
are computed as one plus the maximum fanin level, either incrementally at
construction time (as above) or in bulk by \id{Abc\_NtkLevel}
(\file{src/base/abc/abcDfs.c}).

\section{Object and network lifecycle}

\subsection{Creating and deleting objects}

\id{Abc\_ObjAlloc} (\file{src/base/abc/abcObj.c}) fetches zeroed memory from the
network's fixed-size manager \id{pMmObj}, stamps \id{pNtk} and \id{Type}, and
leaves \id{Id = -1}. \id{Abc\_NtkCreateObj} assigns the ID (its index in
\id{vObjs}), pushes the object onto \id{vObjs}, bumps
\id{nObjCounts[Type]} and \id{nObjs}, and files it into the right role array
(\id{vPis}/\id{vCis} for a PI, \id{vBoxes} for a latch, etc.). The thin
per-type wrappers \id{Abc\_NtkCreatePi}, \id{Abc\_NtkCreateNode},
\id{Abc\_NtkCreateLatch}, \dots\ all funnel through it.

\id{Abc\_NtkDeleteObj} reverses this: it strips the object's name, unhooks every
fanin/fanout edge, sets its \id{vObjs} slot to \id{NULL} (leaving a hole),
decrements the counters, and removes it from its role array.
\id{Abc\_ObjRecycle} then frees the fanin/fanout arrays and returns the memory to
\id{pMmObj}. Because deletion leaves holes, \id{Abc\_NtkObjNum} (live count)
differs from \id{Abc\_NtkObjNumMax} (the size of \id{vObjs}).

\subsection{Duplicating a network}

Two routines in \file{src/base/abc/abcNtk.c} anchor most transformations:

\begin{itemize}
\item \id{Abc\_NtkStartFrom} allocates an empty network of the requested
\id{Type}/\id{Func}, copies the name/spec, clears every \id{pCopy}, and clones
just the CIs, COs, and boxes --- setting each source object's \id{pCopy} to its
new counterpart. The interior is intentionally left empty for the caller to
rebuild.
\item \id{Abc\_NtkDup} calls \id{Abc\_NtkStartFrom} with the same type, then
fills the interior. For a strashed network it walks the ANDs in order and calls
\id{Abc\_AigAnd(\dots, Abc\_ObjChild0Copy(pObj), Abc\_ObjChild1Copy(pObj))},
so the copy is re-hashed as it is built; for a logic/netlist network it dups the
remaining objects and reconnects fanins through \id{pCopy}.
\end{itemize}

\begin{lstlisting}[style=abccode]
// inside Abc_NtkDup, for a strashed network:
Abc_AigForEachAnd( pNtk, pObj, i )
    pObj->pCopy = Abc_AigAnd( (Abc_Aig_t *)pNtkNew->pManFunc,
                              Abc_ObjChild0Copy(pObj), Abc_ObjChild1Copy(pObj) );
Abc_NtkForEachCo( pNtk, pObj, i )
    Abc_ObjAddFanin( pObj->pCopy, Abc_ObjChild0Copy(pObj) );
\end{lstlisting}

This \id{pCopy}-threading idiom --- clear copies, map terminals, rebuild
interior, reconnect outputs --- recurs throughout ABC and is worth recognizing
on sight.
